\begin{frame}{\Zoo en pratique}
\Huge
\centering
\begin{tabular}{ccccc}
      \begin{tabular}{c}
          \includegraphics[scale=0.25]{images/ocaml.png}
        \\
          \includegraphics[scale=0.1]{images/dune.png}
      \end{tabular}
    &&
      $\stackrel{\mbox{\Large\texttt{ocaml2zoo}}}{\longrightarrow}$
    &&
      \begingroup
        \renewcommand{\arraystretch}{1.5}
        \begin{tabular}{c}
            \includegraphics[scale=0.5]{images/iris.png}
          \\
            \textcolor{Green}{\Zoo}
          \\
            \includegraphics[scale=0.25]{images/rocq.png}
        \end{tabular}
      \endgroup
\end{tabular}
\end{frame}

\newcommand{\colorDomainslib}{blue}
\newcommand{\colorSaturn}{red}
\newcommand{\colorScheduler}{cyan}
\newcommand{\colorQueue}{orange}

\begin{frame}[fragile]{\Zoo en pratique}
\centering
\begin{tabular}{ccc}
    \begin{minipage}{0.4\textwidth}\dirtree{%
      .1 project.
        .2 dune-project.
        .2 lib.
          .3 \textcolor{\colorDomainslib}{domainslib}.
            .4 dune.
            .4 \textcolor{\colorScheduler}{scheduler.ml}.
            .4 scheduler.mli.
          .3 \textcolor{\colorSaturn}{saturn}.
            .4 dune.
            .4 \textcolor{\colorQueue}{queue.ml}.
            .4 queue.mli.
    }\end{minipage}
  &
    $\Longrightarrow$
  &
    \begin{minipage}{0.4\textwidth}\dirtree{%
      .1 theories.
        .2 \textcolor{\colorDomainslib}{domainslib}.
          .3 \textcolor{\colorScheduler}{scheduler\_\_code.v}.
          .3 \textcolor{\colorScheduler}{scheduler\_\_types.v}.
        .2 \textcolor{\colorSaturn}{saturn}.
          .3 \textcolor{\colorQueue}{queue\_\_code.v}.
          .3 \textcolor{\colorQueue}{queue\_\_types.v}.
    }\end{minipage}
\end{tabular}
\vfill
\Large
\begin{center}
  \texttt{\$ ocaml2zoo project theories}
\end{center}
%\begin{overbox}<2>
%  \centering
%  \texttt{queue.ml}
%  \medskip
%  \begin{minted}{ocaml}
%  let pop t = ...
%  \end{minted}
%\end{overbox}
%\begin{overbox}<4>
%  \centering
%  \texttt{queue\_\_code.v}
%  \medskip
%  \begin{minted}{coq}
%  Definition queue_pop : val := ...
%  \end{minted}
%\end{overbox}
%\begin{overbox}<6>
%  \centering
%  \texttt{scheduler.ml}
%  \medskip
%  \begin{minted}{ocaml}
%  let get_task =
%    ... Saturn.Queue.pop ...
%  \end{minted}
%\end{overbox}
%\begin{overbox}<8>
%  \centering
%  \texttt{scheduler\_\_code.v}
%  \medskip
%  \begin{minted}{coq}
%  From saturn Require Import
%    queue__code
%    queue__types.
%  
%  Definition scheduler_get_task : val :=
%    ... queue_pop ...
%  \end{minted}
%\end{overbox}
%\only<9>{}
\end{frame}

\begin{frame}[fragile]{\Zoo en pratique}
\centering
\begin{tabular}{ccc}
    \begin{minipage}{0.4\textwidth}\dirtree{%
      .1 project.
        .2 dune-project.
        .2 lib.
          .3 \textcolor{\colorDomainslib}{domainslib}.
            .4 dune.
            .4 \textcolor{\colorScheduler}{scheduler.ml}.
            .4 scheduler.mli.
    }\end{minipage}
  &
    $\Longrightarrow$
  &
    \begin{minipage}{0.4\textwidth}\dirtree{%
      .1 theories.
        .2 \textcolor{\colorDomainslib}{domainslib}.
          .3 \textcolor{\colorScheduler}{scheduler\_\_code.v}.
          .3 \textcolor{\colorScheduler}{scheduler\_\_types.v}.
    }\end{minipage}
\end{tabular}
\vfill
\Large
\begin{center}
  \texttt{\$ ocaml2zoo project theories}
\end{center}
%\begin{overbox}<2>
%  \centering
%  \texttt{scheduler\_\_code.v}
%  \medskip
%  \begin{minted}{coq}
%  From saturn Require Import
%    queue__code
%    queue__types.
%  
%  Definition scheduler_get_task : val :=
%    ... queue_pop ...
%  \end{minted}
%\end{overbox}
%\only<3>{}
\end{frame}

\begin{frame}[fragile]{\Zoo en pratique}
\begin{minipage}{0.49\textwidth}
  \begin{minted}{coq}
Lemma stack_push_spec_seq t ι v :
  {{{
    stack_model t vs
  }}}
    stack_push t v
  {{{
    RET ();
    stack_model t (v :: vs)
  }}}.
Proof.
  ...
Qed.
  \end{minted}
\end{minipage}
\begin{minipage}{0.5\textwidth}
  \begin{minted}{coq}
Lemma stack_push_spec_atomic t ι v :
  <<< 
    stack_inv t ι
  | ∀∀ vs,
    stack_model t vs 
  >>>
    stack_push t v @ ↑ι
  <<<
    stack_model t (v :: vs)
  | RET (); True
  >>>.
Proof.
  ...
Qed.
  \end{minted}
\end{minipage}
\end{frame}

\begin{frame}[fragile]{Types algébriques de données}
\begin{minipage}{0.49\textwidth}
  \begin{minted}{ocaml}
type 'a t =
  | Nil
  | Cons of 'a * 'a t

let rec map fn t =
  match t with
  | Nil -> Nil
  | Cons (x, t) ->
      let y = fn x in
      Cons (y, map fn t)
  \end{minted}
\end{minipage}
\begin{minipage}{0.5\textwidth}
  \begin{minted}{coq}
Notation "'Nil'"  := (
  in_type "t" 0
)(in custom zoo_tag).
Notation "'Cons'" := (
  in_type "t" 1
)(in custom zoo_tag).

Definition map : val :=
  rec: "map" "fn" "t" =>
    match: "t" with
    | Nil => §Nil
    | Cons "x" "t" =>
        let: "y" := "fn" "x" in
        ‘Cons( "y", "map" "fn" "t" )
    end.
  \end{minted}
\end{minipage}
\end{frame}

\begin{frame}[fragile]{Enregistrements}
\begin{minipage}{0.49\textwidth}
  \begin{minted}{ocaml}
type 'a t =
  { mutable f1: 'a;
    mutable f2: 'a;
  }

let swap t =
  let f1 = t.f1 in
  t.f1 <- t.f2 ;
  t.f2 <- f1
  \end{minted}
\end{minipage}
\begin{minipage}{0.5\textwidth}
  \begin{minted}{coq}
Notation "'f1'" := (
  in_type "t" 0
)(in custom zoo_field).
Notation "'f2'" := (
  in_type "t" 1
)(in custom zoo_field).

Definition swap : val :=
  fun: "t" =>
    let: "f1" := "t".{f1} in
    "t" <-{f1} "t".{f2} ;;
    "t" <-{f2} "f1".
  \end{minted}
\end{minipage}
\end{frame}

\begin{frame}[fragile]{Enregistrements anonymes}
\begin{minipage}{0.49\textwidth}
  \begin{minted}{ocaml}
type 'a node =
  | Null
  | Node of
    { mutable next: 'a node;
      mutable data: 'a;
    }
  \end{minted}
\end{minipage}
\begin{minipage}{0.5\textwidth}
  \begin{minted}{coq}
Notation "'Null'" := (
  in_type "node" 0
)(in custom zoo_tag).
Notation "'Node'" := (
  in_type "node" 1
)(in custom zoo_tag).

Notation "'next'" := (
  in_type "node__Node" 0
)(in custom zoo_field).
Notation "'data'" := (
  in_type "node__Node" 1
)(in custom zoo_field).
  \end{minted}
\end{minipage}
\end{frame}

\begin{frame}[fragile]{Fonctions mutuellement récursives}
\begin{minipage}{0.49\textwidth}
  \begin{minted}{ocaml}
let rec f x = g x
and g x = f x
  \end{minted}
\end{minipage}
\begin{minipage}{0.5\textwidth}
  \begin{minted}{coq}
Definition f_g := (
  recs: "f" "x" => "g" "x"
  and:  "g" "x" => "f" "x"
)%zoo_recs.

(* boilerplate *)

Definition f := ValRecs 0 f_g.
Definition g := ValRecs 1 f_g.

Instance : AsValRecs' f 0 f_g [f;g].
Proof. done. Qed.
Instance : AsValRecs' g 1 f_g [f;g].
Proof. done. Qed.
  \end{minted}
\end{minipage}
\end{frame}

\begin{frame}[fragile]{Concurrence}
\begin{adjustwidth}{-1em}{-1em}
\centering
\begin{tabular}{ll}
    \mintinline[escapeinside=||]{ocaml}{Atomic.set |$e_1$| |$e_2$|} &
    $e_1\ \texttt{<-}\ e_2$
  \\
    \mintinline[escapeinside=||]{ocaml}{Atomic.exchange |$e_1$| |$e_2$|} &
    $\texttt{Xchg}\ e_1 \texttt{.[contents]}\ e_2$
  \\
    \mintinline[escapeinside=||]{ocaml}{Atomic.compare_and_set |$e_1$| |$e_2$| |$e_3$|} &
    $\texttt{CAS}\ e_1 \texttt{.[contents]}\ e_2\ e_3$
  \\
    \mintinline[escapeinside=||]{ocaml}{Atomic.fetch_and_add |$e_1$| |$e_2$|} &
    $\texttt{FAA}\ e_1 \texttt{.[contents]}\ e_2$
  \\\\
    \mintinline[escapeinside=||]{ocaml}{type t = { ...; mutable |$f$|: |$\tau$| [@atomic]; ... }}
  \\
    \mintinline[escapeinside=||]{ocaml}{Atomic.Loc.exchange [%atomic.loc |$e_1$|.|$f$|] |$e_2$|} &
    $\texttt{Xchg}\ e_1 \texttt{.[} f \texttt{]}\ e_2$
  \\
    \mintinline[escapeinside=||]{ocaml}{Atomic.Loc.compare_and_set [%atomic.loc |$e_1$|.|$f$|] |$e_2$| |$e_3$|} &
    $\texttt{CAS}\ e_1 \texttt{.[} f \texttt{]}\ e_2\ e_3$
  \\
    \mintinline[escapeinside=||]{ocaml}{Atomic.Loc.fetch_and_add [%atomic.loc |$e_1$|.|$f$|] |$e_2$|} &
    $\texttt{FAA}\ e_1 \texttt{.[} f \texttt{]}\ e_2$
\end{tabular}
\vfill
\url{https://github.com/ocaml/ocaml/pull/13404}
\url{https://github.com/ocaml/ocaml/pull/13707}
\end{adjustwidth}
\end{frame}

\begin{frame}{Bibliothèque standard}
\LARGE
\begin{minipage}{0.5\textwidth}
  \begin{itemize}
    \item \texttt{Array}
    \item \texttt{Dynarray}
    \item \texttt{List}
    \item \texttt{Stack}
    \item \texttt{Queue}
    \item \texttt{Deque}
  \end{itemize}
\end{minipage}
\begin{minipage}{0.4\textwidth}
  \begin{itemize}
    \item \texttt{Domain}
    \item \texttt{Atomic\_array}
    \item \texttt{Mutex}
    \item \texttt{Condition}
  \end{itemize}
\end{minipage}
\end{frame}
